\scanpage{31}
\mspage{30}

\begin{center}$-\ 47\ -$\end{center}\hfill n\textsuperscript{o} 272

$n=[L:K]$.\quad b) \uline{L'intersection} $B$ \uline{des} $A_i$ \uline{est la clôture intégrale de} $A$ dans $L$.

c) \uline{Si l'on pose} $\underline{\mathrm p}_i=B\cap\underline{\mathrm m}(A_i)$, \uline{les idéaux} $\underline{\mathrm p}_i$ \uline{sont des idéaux maximaux de} $B$, \uline{distincts} ; \uline{tout idéal maximal de} $B$ \uline{est égal à l'un des} $\underline{\mathrm p}_i$ \uline{et l'on a}
\[
A_i=B_{\underline{\mathrm p}_i}.
\]
Nous noterons $v_i$ (resp.$v$) des valuations associées aux $A_i$ (resp. à $A$), et $\Gamma_i$ (resp.$\Gamma$) les groupes des ordres correspondants. Le groupe $\Gamma$ s'identifie à un sous-groupe de chacun des groupes $\Gamma_i$.

Montrons d'abord que l'on ne peut pas avoir $A_i\supset A_j$ si $i\ne j$. En effet, d'après ce que l'on a vu au \S4, le groupe $\Gamma_i$ s'identifierait à un groupe quotient $\Gamma_j/H$, où $H$ est un sous-groupe isolé de $\Gamma_j$ ; puisque l'application canonique $\Gamma\longrightarrow\Gamma_j\longrightarrow\Gamma_j/H=\Gamma_i$ est une injection, on a $\Gamma\cap H=\{0\}$ dans $\Gamma_j$, ce qui montre que $H$ est isomorphe à un sous-groupe de $\Gamma_j/\Gamma$ ; comme ce dernier groupe est fini (prop.1), on a $H=0$, et les valuations $v_i$ et $v_j$ sont équivalentes, d'où $A_i=A_j$ contrairement à l'hypothèse $i\ne j$.

Montrons maintenant que le nombre des $A_i$ est \unsure{'} $\leqslant n$. Supposons donc qu'il existe $s$ anneaux $A_i$ distincts, avec $s>n$. D'après le cor. 2 à la prop.2 du \S6, (qui est applicable, en vertu de ce que l'on vient de voir), il existe des éléments $x_i\in L$, $1\leqslant i\leqslant s$, tels que $v_i(x_i)=0$ et $v_j(x_i)>0$ pour $i\ne j$. Montrons que les $x_i$ sont linéairement indépendants sur $K$, ce qui fournira la contradiction cherchée. Supposons que l'on ait une relation linéaire non triviale :
\[
\sum a_ix_i=0,\text{ avec }a_i\in K.
\]
Soit $a_j$ celui des $a_i$ dont la valuation est minimum ; quitte à multiplier les $a_i$ par $a_j^{-1}$ on peut supposer que $a_j=1$. On a alors une relation :
